% Standalone problem document, generated from the adjacent README.md.
% Compile with XeLaTeX. Mathematical content is maintained in Markdown.
\documentclass[11pt,a4paper]{article}
\usepackage[fontsize=10.5pt]{scrextend}
\usepackage[margin=22mm,headheight=15pt,headsep=9mm,footskip=11mm]{geometry}
\usepackage{fontspec}
\setmainfont{texgyrepagella}[Extension=.otf,UprightFont=*-regular,BoldFont=*-bold,ItalicFont=*-italic,BoldItalicFont=*-bolditalic]
\setsansfont{texgyreheros}[Extension=.otf,UprightFont=*-regular,BoldFont=*-bold,ItalicFont=*-italic,BoldItalicFont=*-bolditalic]
\setmonofont{texgyrecursor}[Extension=.otf,UprightFont=*-regular,BoldFont=*-bold,ItalicFont=*-italic,BoldItalicFont=*-bolditalic,Scale=MatchLowercase]
\usepackage{amsmath,amssymb}
\usepackage{unicode-math}
\setmathfont{texgyrepagella-math.otf}
\usepackage{xcolor,microtype,enumitem,needspace,fancyhdr}
\usepackage{bookmark}
\definecolor{ink}{HTML}{17344B}
\definecolor{accent}{HTML}{197D80}
\definecolor{muted}{HTML}{596773}
\hypersetup{colorlinks=true,linkcolor=accent,urlcolor=accent,pdftitle={MI-32 - Spectral
norms of independent entries with regular moment
growth},pdfauthor={Open Problems in Numerical Linear Algebra}}
\urlstyle{same}
\setlength{\parindent}{0pt}
\setlength{\parskip}{5pt plus 1pt minus 1pt}
\linespread{1.04}
\setlength{\emergencystretch}{3em}
\widowpenalty=10000
\clubpenalty=10000
\displaywidowpenalty=10000
\allowdisplaybreaks[1]
\setlist{nosep,leftmargin=1.5em}
\providecommand{\tightlist}{\setlength{\itemsep}{0pt}\setlength{\parskip}{0pt}}
\providecommand{\pandocbounded}[1]{#1}
\pagestyle{fancy}
\fancyhf{}
\fancyhead[L]{\sffamily\footnotesize\color{muted}OPEN PROBLEMS IN NUMERICAL LINEAR ALGEBRA}
\fancyhead[R]{\sffamily\bfseries\footnotesize\color{accent}MI-32}
\fancyfoot[L]{\parbox[b]{0.9\textwidth}{\sffamily\fontsize{8}{10}\selectfont\color{muted}Editorial ratings; a bounded search cannot certify that a problem remains unsolved.\\Literature check: 2026-09-14}}
\fancyfoot[R]{\sffamily\footnotesize\color{muted}\thepage}
\renewcommand{\headrulewidth}{0.3pt}
\renewcommand{\headrule}{\hbox to\headwidth{\color{accent}\leaders\hrule height \headrulewidth\hfill}}
\makeatletter
\renewcommand{\section}{\Needspace{5\baselineskip}\@startsection{section}{1}{0pt}{12pt plus 2pt minus 2pt}{4pt}{\sffamily\bfseries\large\color{ink}}}
\renewcommand{\subsection}{\Needspace{4\baselineskip}\@startsection{subsection}{2}{0pt}{9pt plus 1pt minus 1pt}{3pt}{\sffamily\bfseries\normalsize\color{ink}}}
\makeatother
\setcounter{secnumdepth}{0}
\begin{document}
{\sffamily\small\color{muted}Matrix inequalities and norms\par}
\vspace{3pt}
{\raggedright\hyphenpenalty=10000\exhyphenpenalty=10000\sffamily\bfseries\fontsize{21}{24}\selectfont\color{ink}Spectral
norms of independent entries with regular moment growth\par}
\vspace{7pt}
{\sffamily\small\textbf{Difficulty:} challenging\quad\textbf{Importance:} interesting
to the community\par}
{\sffamily\small\bfseries\color{accent}Status: Lean verified\par}
\vspace{4pt}
{\color{accent}\hrule height 0.8pt}
\vspace{6pt}
\textbf{Topic:} Expected spectral norms of inhomogeneous random matrices

\textbf{Lean verified --- 2026-09-14. Formalization: Diar Heidary,
unaffiliated.} The upper comparison displayed below, which is the whole
of the conjecture stated in this entry, is proved in Lean 4 and
machine-checked. The proof source, its mechanical verification and its
transitive axiom report are recorded in
\href{https://github.com/ajt60gaibb/OpenProblemsInNLA/blob/main/matrix-inequalities-and-norms/MI-32\#lean-proof-and-verification-evidence--2026-09-14}{Lean
proof and verification evidence} below. The conjecture remains credited
to Latała and Świątkowski, and the reverse inequality quoted in this
entry remains their Theorem 4.1 and is not formalized. AI assistance is
disclosed; no external human peer review is claimed.

\textbf{Rating rationale:} The missing upper estimate must account for
localized random fluctuations across all variance profiles. Even
weighted random signs retain an iterated-logarithm gap. A sharp estimate
would improve quantitative understanding of the largest singular value
in random matrix models.

\subsection{Statement}\label{statement}

For a real random variable \(Z\) and \(r>0\), write
\(\|Z\|_{L_r}=(\mathbb E|Z|^r)^{1/r}\). Fix \(\alpha\ge1\). Let
\(X=(X_{ij})\in\mathbb R^{n\times n}\) have independent mean-zero
entries with finite absolute moments of every order, satisfying

\[\|X_{ij}\|_{L_{2r}}\le\alpha\|X_{ij}\|_{L_r}
\qquad(1\le i,j\le n,\ r\ge1).\]

Set \([n]=\{1,\ldots,n\}\) and

\[M(X)=\max_i\left(\sum_j\mathbb E X_{ij}^2\right)^{1/2}
+\max_j\left(\sum_i\mathbb E X_{ij}^2\right)^{1/2},\]

\[D(X)=\max_{1\le k\le n}\;
\min_{\substack{I\subseteq[n]\\|I|\le k}}
\sup_{\substack{s,t\in\mathbb R^n\\\|s\|_2,\|t\|_2\le1}}
\left\|\sum_{\substack{i\notin I\\j\notin I}}X_{ij}s_it_j
\right\|_{L_{\ln(k+1)}}.\]

\textbf{Conjecture (Latała--Świątkowski).} For every \(\alpha\ge1\)
there is a constant \(C_\alpha>0\) such that, for every \(n\ge1\) and
every such matrix,

\[\mathbb E\|X\|_2\le C_\alpha\bigl(M(X)+D(X)\bigr),\]

where \(\|X\|_2\) is the spectral norm. The constant is independent of
the dimension and entry laws. The same deterministic index set \(I\)
deletes both rows and columns; the supremum is outside the
random-variable moment. The definition includes \(L_{\ln2}\) when
\(k=1\), using the displayed moment functional even though its exponent
is below one.

The reverse inequality up to a constant depending only on \(\alpha\) is
proved, so the target is equivalent to the source's two-sided
comparison. No identical-distribution or symmetry assumption is imposed
on the entries.

\subsection{Lean proof and verification evidence ---
2026-09-14}\label{lean-proof-and-verification-evidence-2026-09-14}

\subsubsection{Proof source, toolchain and pinned
dependencies}\label{proof-source-toolchain-and-pinned-dependencies}

The formalization is the Lean 4 project
\href{https://github.com/DiarHaidary/Spectral-norms-of-independent-entries-with-regular-moment-growth}{DiarHaidary/Spectral-norms-of-independent-entries-with-regular-moment-growth}
at the immutable revision
\href{https://github.com/DiarHaidary/Spectral-norms-of-independent-entries-with-regular-moment-growth/tree/762bd5ec5050a96f5e6ba3926b6cda4816fcd4b0}{\textit{762bd5ec5050a96f5e6ba3926b6cda4816fcd4b0}}.
That exact commit is registered with the Palomar registry as
\href{https://palomar-registry.org/entry.html?id=PALOMAR-2026-09-14-000006&version=1}{PALOMAR-2026-09-14-000006,
version 1}, which also preserves an archival fork of it.

\begin{itemize}
\tightlist
\item
  Lean toolchain: \textit{leanprover/lean4:v4.33.0}.
\item
  \textit{mathlib} pinned at
  \textit{db584cd6d46c92f209a44c0f1c829460d327499d}; the remaining eight
  Lake dependencies are pinned to exact commits in the committed
  \textit{lake-manifest.json}.
\item
  \textit{Challenge.lean} imports Mathlib only.
\end{itemize}

\subsubsection{Compared declaration and its correspondence to this
entry}\label{compared-declaration-and-its-correspondence-to-this-entry}

Comparator selects the single declaration \textit{MI32.main\_upper},
with \textit{definition\_names} empty, so no part of the statement is a
replaceable hole. Its literal Lean text is in
\href{https://github.com/DiarHaidary/Spectral-norms-of-independent-entries-with-regular-moment-growth/blob/762bd5ec5050a96f5e6ba3926b6cda4816fcd4b0/Challenge.lean}{\textit{Challenge.lean}},
the file a reader is expected to audit. In the notation of this entry it
asserts

\[\forall\,\alpha\ge1\quad\exists\,C>0\quad\text{such that}\quad
\mathbb E\|X\|_2\le C\bigl(M(X)+D(X)\bigr)\]

for every \(n\ge1\) and every admissible \(X\). The correspondence
between the displayed quantities and the definitions in that file is:

\begin{itemize}
\tightlist
\item
  \(\|Z\|_{L_r}\) is \textit{moment}, defined as \((\int|Z|^r)^{1/r}\)
  with a \emph{real} exponent. It is therefore the displayed functional
  also when \(r<1\), so the order \(\ln2\) at \(k=1\) is kept literally
  and is not replaced by one.
\item
  The hypotheses on \(X\) are \textit{RegularEntries}: entry
  measurability; joint independence of the family indexed by ordered
  pairs, as \textit{iIndepFun}; each entry has Bochner integral zero;
  \(|X_{ij}|^p\) is integrable for every real \(p>0\); and
  \(\|X_{ij}\|_{L_{2r}}\le\alpha\|X_{ij}\|_{L_r}\) for every real
  \(r\ge1\). Neither identical distribution nor symmetry is assumed.
\item
  \(\|X\|_2\) is \textit{spectralNorm}, the operator norm of the
  Euclidean linear map of the matrix.
\item
  \(M(X)\) is \textit{varianceScale}, the supremum of the row standard
  deviations plus the supremum of the column standard deviations.
\item
  The inner supremum over \(\|s\|_2,\|t\|_2\le1\) after deleting \(I\)
  is \textit{weakMoment}: a real supremum, over deterministic test
  vectors, of the scalar moment of \textit{deletedBilinear}. That
  bilinear form sums over \textit{Finset.univ \textbackslash{}\ I} on
  both axes, so one and the same deterministic set deletes rows and
  columns, and the supremum sits outside the random-variable moment.
\item
  \(D(X)\) is \textit{deletionScale}, the supremum over \(1\le k\le n\)
  of the infimum, over index sets of cardinality at most \(k\), of
  \textit{weakMoment} at the order \(\ln(k+1)\).
\item
  The conjectured inequality is \textit{UpperBoundAt}, quantified over
  every \(n\ge1\), every type, every measurable space on it, every
  probability measure and every such \(X\). The constant is bound before
  all of those, so it depends only on \(\alpha\) and not on the
  dimension or the entry laws.
\end{itemize}

The probability space is an arbitrary type rather than a finite one, and
no estimate enters as a hypothesis or as a custom axiom. The constant is
explicit: \textit{SymmetricDeletion.deletionConstant} evaluated at
\(2\alpha\), times \(e\). It is conservative and is not claimed to be
sharp.

The hypothesis class is not vacuous.
\textit{MI32.RegularWitness.exists\_regularEntries} exhibits an
independent fair-sign matrix satisfying the hypotheses at \(\alpha=1\)
with \(M(X)=\sqrt n>0\), so the displayed bound is a genuine constraint
rather than a statement about an empty class.

\subsubsection{Verification record}\label{verification-record}

This catalog \textbf{reviewed a public verification record}; it did not
rerun the checks itself. Palomar's mechanical verification of the
registered commit ran on Palomar's own Linux infrastructure on
2026-09-14T14:07:31Z:

\begin{itemize}
\tightlist
\item
  \href{https://github.com/PalomarRegistry/PalomarSubmission/actions/runs/34852526384}{workflow
  run 34852526384} in \textit{PalomarRegistry/PalomarSubmission};
\item
  Comparator at \textit{575674928e239f5bc452aab72d1dd7b0f1326494},
  checking that the Solution declaration proves the Challenge statement
  in a separate environment;
\item
  an independent NanoDa kernel replay at
  \textit{68d5ca9db226849b41a6fff59d796ff19d0a8840};
\item
  \textit{lean4export} at
  \textit{15f6055e299ad5b89345e533cc2192f4cc00f659} and the
  \textit{landrun} sandbox at
  \textit{811cfff51ceaf3d9843708aa6d22e9b84ccac8b4}.
\end{itemize}

The registry entry records trust level \textit{high}, and its AI
editorial review returned outcome \textit{neutral}, meaning that no
blocking problem was identified.

The author's own local checks are reproducible from the repository root
with:

\begin{verbatim}
lake exe cache get
lake build
lake env lean Audit.lean
lake build Challenge Solution
lake env lean FullTargetAudit.lean
python3 scripts/verify.py
\end{verbatim}

The last of these exits zero only when the whole gate passes, and writes
a dated record to \textit{verification/status.json}. A successful build
is not on its own evidence of the result; the Comparator and
kernel-replay run above is.

\subsubsection{Transitive axiom report}\label{transitive-axiom-report}

\begin{verbatim}
'MI32.main_upper' depends on axioms: [propext, Classical.choice, Quot.sound]
\end{verbatim}

produced by \textit{lake env lean FullTargetAudit.lean} and retained
at \textit{verification/full-target-axioms.log}. The Comparator
configuration permits only those three axioms and the mechanical run
enforced that. There is no \textit{sorryAx}, no custom axiom and no
reliance on native execution. \textit{Challenge.lean} retains a
deliberate \textit{sorry}, which is the Challenge hole by construction
and is never imported by the proof.

\subsubsection{Scope, attribution and
limits}\label{scope-attribution-and-limits}

\begin{itemize}
\tightlist
\item
  \textbf{Scope.} Only the upper comparison displayed in this entry is
  formalized. The reverse inequality, Theorem 4.1 of
  Latała--Świątkowski, is quoted in this entry for context and is not
  formalized or selected. No quantile or two-sided statement is claimed,
  and no claim is made about the sharpness of the constant's dependence
  on \(\alpha\).
\item
  \textbf{Mathematical attribution.} The conjecture is Conjecture 4.3 of
  Latała--Świątkowski. The proof argument and its formalization are by
  Diar Heidary. The manuscript is
  \href{https://github.com/DiarHaidary/Spectral-norms-of-independent-entries-with-regular-moment-growth/blob/bd00df8dca9c6128cf528591d3485f8ca9eecd88/docs/manuscript/MI32.pdf}{MI32.pdf},
  with its
  \href{https://github.com/DiarHaidary/Spectral-norms-of-independent-entries-with-regular-moment-growth/blob/bd00df8dca9c6128cf528591d3485f8ca9eecd88/docs/manuscript/MI32.tex}{LaTeX
  source}; its formal-verification section tabulates, statement by
  statement, which Lean declaration carries each numbered result.
\item
  \textbf{Deliberate divergences from the manuscript}, recorded in the
  project's \textit{docs/SOURCE\_FIDELITY.md}: the formal far-remainder
  estimate avoids the variance-envelope input of Latała--van
  Handel--Youssef by raising the screening threshold and using an
  elementary centered-Gram second moment, and the formal decomposition
  runs on the original rectangular matrix rather than on the
  matrix-symmetric dilation.
\item
  \textbf{AI assistance is disclosed.} The proof, the Lean development
  and the exposition were produced with substantial AI-agent assistance
  under the author's direction. The statement-fidelity review and the
  registry's editorial review are AI reviews. Neither is human peer
  review and none is claimed. Registration with Palomar is not
  acceptance, endorsement or a novelty claim.
\end{itemize}

\subsection{Known cases and numerical
significance}\label{known-cases-and-numerical-significance}

The source proves the Gaussian-mixture case under the same moment
condition. Weighted signs \(X_{ij}=a_{ij}\varepsilon_{ij}\), with
independent fair \(\varepsilon_{ij}\in\{-1,1\}\), satisfy the condition
with \(\alpha=1\) and form the source's earlier Conjecture 1.2. Latała
proves this case when \(a_{ij}\in\{0,1\}\) and proves the general
weighted-sign estimate with an additional factor of order
\(\log\log\log n\). Meller's 2026 result gives an iterated-logarithm
loss for a wider class of symmetric entries; it does not remove that
loss. These cases remain grouped in this entry.

The spectral norm is the largest singular value and measures the largest
Euclidean amplification by a random matrix. The formula combines row and
column variance scales with moments of bilinear forms after deleting a
limited number of coordinates, accounting for concentration on small
parts of a matrix that simpler variance-only bounds can miss.

\subsection{References and status
check}\label{references-and-status-check}

\begin{itemize}
\tightlist
\item
  R. Latała and W. Świątkowski, \emph{Norms of randomized circulant
  matrices}, Electronic Journal of Probability 27 (2022), paper 80,
  1--23. \href{https://doi.org/10.1214/22-EJP799}{DOI};
  \href{https://arxiv.org/pdf/2106.03139v2}{current arXiv v2}, dated
  2022-05-27. Conjecture 4.3 on preprint p.25, with condition (25) on
  p.24, states the target. Theorem 4.1 supplies the lower bound;
  Proposition 4.4 and the paragraph before it cover Gaussian mixtures.
  Conjecture 1.2 on p.2 is the weighted-sign special case.
\item
  R. Latała, \emph{On the spectral norm of Rademacher matrices}.
  \href{https://doi.org/10.1090/tran/9637}{DOI};
  \href{https://arxiv.org/html/2405.13656v2}{current arXiv v2}, dated
  2025-08-18. Equations (1.2)--(1.3), Theorem 1.1, and Theorem 1.9 give
  the weighted-sign conjecture and the stated partial results. This
  paper uses the comparable truncated logarithm \(\max\{1,\ln k\}\) in
  place of \(\ln(k+1)\).
\item
  R. Meller, \emph{Spectral norm of matrices with independent entries up
  to polyloglog},
  \href{https://arxiv.org/html/2512.23673v2}{arXiv:2512.23673v2}, dated
  2026-01-29, introduction, equation (3), and Theorem 1.1. It explicitly
  identifies the remaining logarithmic gap and Conjecture 4.3.
\end{itemize}

On 2026-09-11, checked the original paper's current version, both later
papers, and exact-title, author/conjecture, Rademacher spectral-norm,
proof/counterexample and 2026 searches. No full resolution was found.
Gaussian operator-norm results and the general Bernoulli-process theorem
do not prove this displayed formula. This is a bounded literature check.
\end{document}
